\documentclass[10pt,a4paper]{article}

\usepackage[margin=1in]{geometry}
\usepackage{amsmath,amssymb,amsfonts,amsthm}
\usepackage{graphicx}
\usepackage{booktabs}
\usepackage{hyperref}
\usepackage{xcolor}
\usepackage{microtype}
\usepackage{multirow}
\usepackage{array}
\usepackage{enumitem}
\usepackage{titlesec}
\usepackage{abstract}

\hypersetup{
    colorlinks=true,
    linkcolor=blue,
    citecolor=blue,
    urlcolor=blue
}

\title{\textbf{TensionLM: Sigmoid Tension as Constraint Relaxation\\for Language Modelling}}

\author{
    BoggersTheFish\\
    Independent Researcher\\
    \texttt{\href{https://github.com/BoggersTheFish/bozo}{github.com/BoggersTheFish/bozo}}
    \quad
    \texttt{\href{https://huggingface.co/BoggersTheFish/TensionLM-117M-Curriculum}{huggingface.co/BoggersTheFish}}
}

\date{April 2026}

\begin{document}

\maketitle

\begin{abstract}
We present \textbf{TensionLM}, a language model architecture motivated by Thinking System (TS) theory, which models all computation as constraint relaxation over a graph: nodes are states, edges are constraints, and $\tau$ (tension) is the degree to which a constraint is active but not yet satisfied. Under this framework, softmax attention is theoretically incorrect --- it is zero-sum, forcing tokens to compete for a fixed probability budget, whereas real constraints operate independently. We replace softmax with \textbf{sigmoid tension attention}, scoring each token pair independently with no global normalisation and normalising outputs by tau-mass ($\Sigma\tau$) rather than position count. At 1.1M parameters, TensionLM matches transformer perplexity exactly (57.7 vs 57.8 on WikiText-2). We introduce four TS-motivated auxiliary losses enforcing transitivity, entropy, manifold closure, and head diversity directly in the constraint graph. Training 117M parameters on a Logic $\rightarrow$ Language $\rightarrow$ Mathematics curriculum achieves first-contact mathematics perplexity of \textbf{24} --- 96$\times$ better than cold start. The tension field demonstrates simultaneous non-competitive constraints impossible under softmax, spontaneous head specialisation without supervision, and 43.5\% accuracy on formal reasoning with no instruction tuning.
\end{abstract}

\section{Introduction}

The transformer \cite{vaswani2017attention} has dominated language modelling for nearly a decade. Its central operation --- softmax attention --- has remained structurally unchanged despite continuous architectural refinement. Softmax computes a probability distribution over context positions, with weights summing to one. This is computationally convenient and empirically effective. It also encodes a specific and arguably incorrect prior about how language works: attention is a competition. If token $A$ scores highly, tokens $B$ and $C$ are forced lower. The attention budget is fixed and must be divided.

Consider the sentence: \textit{``If all mammals are warm-blooded and all whales are mammals, then\ldots''} The correct output depends simultaneously on both \textit{mammals} and \textit{whales} at full strength. These are not competing predecessors. They are simultaneously active logical constraints, and the output requires both of them together. Softmax cannot represent this. Under softmax, the model must apportion its fixed attention budget between the two, weakening at least one. Under a constraint graph view, both are active at independent full strength, and the output is a function of both.

This observation is the empirical foundation of \textbf{Thinking System (TS) theory}, a framework developed by the author proposing that computation, cognition, and information processing are unified under constraint relaxation dynamics. The framework proposes five primitives: nodes (states), edges (constraints), $\tau$ (unresolved tension), computation (finding low-tension equilibrium), and learning (accumulating edges from experience). Representation, under TS, is not a vector --- it is a constraint pattern. A concept is defined by what it constrains, not by what it is.

Under this framework, softmax attention is not merely suboptimal. It implements the wrong mathematical object. Softmax encodes competition. TS requires independence. The natural and theoretically motivated replacement is sigmoid: each constraint is evaluated on its own merits, producing a score in $[0, 1]$ without reference to any other score. The resulting tensor of $\tau$ values is not a side-effect of the forward pass --- it is a directly readable constraint graph showing which tokens are constraining which outputs and at what strength.

This paper presents \textbf{TensionLM}, the implementation of this claim. We train models ranging from 1.1M to 117M parameters across five experiments, and we make the following contributions:

\begin{enumerate}[label=(\arabic*)]
    \item A formally motivated replacement for softmax attention: \textbf{sigmoid tension attention} with tau-mass normalisation, grounded in TS theory.
    \item Four \textbf{TS-native auxiliary losses} that enforce structural coherence directly in the constraint graph during training: ConstraintConsistency, TensionEntropy, ManifoldClosure, and TensionDiversity.
    \item A \textbf{three-stage curriculum} (Logic $\rightarrow$ Language $\rightarrow$ Mathematics) that produces a 96$\times$ improvement in first-contact mathematical perplexity over cold start at 117M scale, with evidence that the effect amplifies with model capacity.
    \item \textbf{Interpretability results} demonstrating that the tension field encodes simultaneous non-competitive constraints as predicted by theory, with spontaneous head specialisation into syntactic, semantic, and diffuse roles without any explicit supervision.
    \item A \textbf{formal reasoning evaluation} on a 23-question benchmark yielding 43.5\% accuracy with no instruction tuning or RLHF, and a key finding that the best reasoning checkpoint precedes full training completion.
\end{enumerate}

The code is available at \url{https://github.com/BoggersTheFish/bozo} and models at \url{https://huggingface.co/BoggersTheFish}.

\section{Background and Related Work}

\subsection{Softmax Attention and Its Properties}

The transformer architecture \cite{vaswani2017attention} defines attention as:
\begin{equation}
    \text{Attention}(Q, K, V) = \text{softmax}\!\left(\frac{QK^\top}{\sqrt{d_k}}\right) V
\end{equation}

The softmax function produces a probability simplex: attention weights are non-negative and sum to one. This property is the source of the competitive prior we identify. Several works have noted pathological behaviours arising from softmax's normalisation, including attention sink phenomena \cite{xiao2023efficient} and the tendency for attention to collapse onto a small number of tokens. Recent work has explored relaxing the normalisation constraint in various ways. \citet{ramapuram2024sigmoid} propose sigmoid attention for vision transformers and find competitive performance with softmax, motivating the mechanism from an energy-based perspective. Our motivation is distinct: we adopt sigmoid not for performance or numerical reasons, but because it is the theoretically correct operation under TS for representing independent constraint activation.

\subsection{Curriculum Learning}

Curriculum learning \cite{bengio2009curriculum} organises training data from simple to complex, mimicking human learning. The approach has been applied to language modelling in various forms. \citet{azerbayev2023llemma} demonstrate that training order and domain composition significantly affect mathematical reasoning in large language models, and that pre-training on mathematical text provides strong transfer. \citet{lewkowycz2022solving} show that continued pre-training on mathematical and scientific text dramatically improves quantitative reasoning. Our curriculum is motivated by a specific TS prediction: formal domains where contradictions are structurally impossible produce coherent constraint graphs, while general web text (which contains abundant contradictory claims) produces constraint loops that never resolve, resulting in fuzzy averaged representations that are neither true nor false.

\subsection{Interpretable Attention}

Mechanistic interpretability research \cite{elhage2021mathematical, olsson2022context} has sought to understand what transformer attention circuits compute, primarily through post-hoc reverse engineering. Probing classifiers, activation patching, and circuit analysis are standard tools. These techniques exist precisely because softmax attention was not designed for interpretability. TensionLM takes the opposite approach: interpretability is an architectural property, not an analysis technique. The $\tau$ tensor produced by the forward pass is a directly readable constraint graph, requiring no reverse engineering.

\subsection{Formal Reasoning in Language Models}

Mathematical and formal reasoning has become a key evaluation domain for large language models \cite{hendrycks2021measuring, cobbe2021training}. Models trained primarily on general web text struggle with multi-step logical inference, and substantial work has gone into chain-of-thought prompting \cite{wei2022chain}, scratchpad methods, and specialised fine-tuning to improve performance. TensionLM takes a pre-training approach: we hypothesise that a constraint graph pre-loaded with logical structure --- through curriculum training on synthetic inference data --- will acquire formal reasoning capabilities that do not require post-hoc instruction tuning.

\subsection{Thinking System Theory}

TS theory is a framework developed independently by the author. It is not the first framework to propose constraint-based or energy-based views of computation --- related ideas appear in Hopfield networks \cite{hopfield1982neural}, energy-based models \cite{lecun2006tutorial}, and the Free Energy Principle in neuroscience \cite{friston2010free}. TS differs in framing: it treats the constraint graph not as an internal implementation detail but as the primary object of interest, and proposes that the correct architecture for a language model should expose this graph directly rather than producing it as a hidden byproduct of softmax normalisation.

\section{Architecture}

\subsection{Sigmoid Tension Attention}

The core architectural change in TensionLM is the replacement of softmax attention with sigmoid tension. For each output position $t$ and window offset $w$, we compute:

\begin{equation}
    \tau[t, w] = \sigma\!\left(\frac{\mathbf{q}_t \cdot \mathbf{k}_{t-w}}{\sqrt{d_{\text{head}}}}\right)
    \label{eq:tau}
\end{equation}

where $\sigma$ denotes the sigmoid function and the window covers $w \in \{1, \ldots, W\}$ past positions causally. The output for position $t$ is then:

\begin{equation}
    \mathbf{o}[t] = \frac{\displaystyle\sum_{w=1}^{W} \tau[t,w] \cdot \mathbf{v}_{t-w}}{\displaystyle\sum_{w=1}^{W} \tau[t,w]}
    \label{eq:output}
\end{equation}

The denominator in Equation~\ref{eq:output} is \textbf{tau-mass normalisation}: we divide by the sum of actual tension values rather than the number of valid positions. This is the TS-correct normalisation. Dividing by position count implicitly assumes each position contributes equally; dividing by $\Sigma\tau$ means the output is proportional to actual constraint strength. If all constraints are weak ($\tau \approx 0$ everywhere), the output approaches zero. If one constraint is strong, it dominates proportionally. This is the correct behaviour for a constraint relaxation system.

The key properties that distinguish Equations~\ref{eq:tau} and~\ref{eq:output} from softmax attention are: (1) there is no global normalisation --- the score at one position does not affect any other; (2) multiple positions can simultaneously hold $\tau \approx 1.0$, representing full simultaneous constraint activation; (3) the $\tau$ tensor is directly interpretable as a constraint graph, with values in $[0, 1]$ representing activation strength.

\subsection{Global Tension Layers}

Local windowed attention (window size $W = 64$) provides computational efficiency but limits long-range constraint propagation. To address this, every fourth TensionBlock uses a global window ($W = 1024$), enabling the constraint graph to encode long-range dependencies without incurring $O(T^2)$ cost at every layer. This hybrid design is motivated by the TS observation that constraints operate at multiple scales: local syntactic constraints within phrases, and global semantic constraints across sentences.

\subsection{Full Architecture}

The full model follows a standard transformer backbone with TensionBlocks replacing standard attention blocks:

\begin{center}
    Embedding $\rightarrow$ $N \times$ TensionBlock $\rightarrow$ RMSNorm $\rightarrow$ LM Head (weight-tied)
\end{center}

Each TensionBlock consists of: RMSNorm $\rightarrow$ MultiHeadCausalTensionLayer $\rightarrow$ residual $\rightarrow$ RMSNorm $\rightarrow$ SwiGLU FFN $\rightarrow$ residual.

\paragraph{Positional encoding.} Rotary position embeddings (RoPE) \cite{su2021roformer} are applied to queries and keys, providing better length generalisation than learned absolute positional embeddings.

\paragraph{Feed-forward network.} We use SwiGLU \cite{shazeer2020glu}: $\text{out} = \text{proj}(\text{silu}(\text{gate}(\mathbf{x})) \odot \text{val}(\mathbf{x}))$, as used in LLaMA and PaLM.

\paragraph{Weight initialisation.} Projection weights are initialised with standard deviation $\sigma / \sqrt{\text{depth}}$, and output projections with $\sigma / \sqrt{2 \cdot \text{depth}}$, ensuring stable gradient flow at 12 layers without learning rate tuning.

\paragraph{Fused Triton kernel.} The forward and backward pass of the tension attention layer are implemented in a single fused Triton kernel. The naive PyTorch implementation requires materialising a $B \times T \times H \times W \times d_h$ intermediate tensor; the fused kernel avoids this entirely, achieving a 64$\times$ memory reduction.

\subsection{Model Configurations}

\begin{table}[h]
\centering
\caption{Model configurations used across experiments.}
\label{tab:configs}
\vspace{4pt}
\begin{tabular}{lrrrrrr}
\toprule
Config & Params & $d_{\text{model}}$ & Layers & Heads & $W$ & Vocab \\
\midrule
Tiny & 1.1M & 128 & 4 & 4 & 32 & 2,048 \\
Small & 13.5M & 256 & 6 & 4 & 32 & 32,768 \\
Large & 117M & 768 & 12 & 12 & 64 & 32,768 \\
\bottomrule
\end{tabular}
\end{table}

\section{TS-Native Training Objectives}

Beyond standard cross-entropy next-token prediction, we introduce four auxiliary losses that enforce TS-predicted structural properties directly in the constraint graph during training. These losses do not supervise the model's outputs --- they supervise the internal structure of the $\tau$ field.

\subsection{Constraint Consistency Loss}

TS predicts that constraint graphs must satisfy transitivity: if concept $A$ constrains concept $B$, and $B$ constrains $C$, then $A$ must also constrain $C$. Transitivity is a basic property of coherent reasoning, and its absence in a constraint graph would indicate that the model's internal representations are locally consistent but globally incoherent. We enforce this directly in the $\tau$ field:

\begin{equation}
    \mathcal{L}_{\text{consistency}} = \mathbb{E}_{t,A,B,C}\!\left[\max\!\left(0,\ \tau[t, A] \cdot \tau[t, B] - \tau[t, C] - \epsilon\right)\right]
\end{equation}

where $A$, $B$, $C$ are token positions forming a potential transitivity chain and $\epsilon$ is a margin. The loss penalises cases where both $A \rightarrow B$ and $B \rightarrow C$ are active ($\tau[t,A]$ and $\tau[t,B]$ are high) but $A \rightarrow C$ is not. Weight: 0.1.

\subsection{Tension Entropy Loss}

The tension entropy loss prevents two degenerate configurations: isolated nodes ($\tau \approx 0$ everywhere, meaning no constraints are active) and saturated nodes ($\tau \approx 1$ everywhere, meaning all constraints are equally active and the graph encodes no structure). Both degeneracies represent failures of the constraint graph to encode useful information. We penalise distributions that deviate from a healthy spread:

\begin{equation}
    \mathcal{L}_{\text{entropy}} = \mathbb{E}_{t}\!\left[\text{ReLU}\!\left(\log W - H(\boldsymbol{\tau}_t)\right)\right]
\end{equation}

where $H(\boldsymbol{\tau}_t) = -\sum_w p_w \log p_w$ with $p_w = \tau[t,w] / \sum_w \tau[t,w]$. The loss fires when the entropy of the normalised tension distribution falls below $\log W$, i.e., when the distribution is too peaked. Weight: 0.05.

\subsection{Manifold Closure Loss}

TS predicts that the constraint field should close on itself across a sequence: the hidden state at the end of a sequence should remain coherent with the state at the beginning. This is analogous to the requirement in physics that a field be consistent across a closed boundary. We enforce this with a cosine similarity loss between the first and last hidden states:

\begin{equation}
    \mathcal{L}_{\text{manifold}} = 1 - \cos(\mathbf{h}_0, \mathbf{h}_{T-1})
\end{equation}

Weight: 0.01.

\subsection{Tension Diversity Loss}

Without explicit encouragement, attention heads can collapse onto a single dominant position, reducing the effective capacity of the multi-head mechanism. The tension diversity loss encourages each head to spread its tension across the window, ensuring that different heads specialise on different aspects of the constraint graph:

\begin{equation}
    \mathcal{L}_{\text{diversity}} = -\mathbb{E}_{\text{heads}}\!\left[H\!\left(\bar{\boldsymbol{\tau}}_{\text{head}}\right)\right]
\end{equation}

where $\bar{\boldsymbol{\tau}}_{\text{head}}$ is the mean tension per position across the sequence for each head. Weight: 0.02.

\subsection{Combined Objective}

The total training loss is:

\begin{equation}
    \mathcal{L} = \mathcal{L}_{\text{CE}} + 0.1 \cdot \mathcal{L}_{\text{consistency}} + 0.05 \cdot \mathcal{L}_{\text{entropy}} + 0.01 \cdot \mathcal{L}_{\text{manifold}} + 0.02 \cdot \mathcal{L}_{\text{diversity}}
\end{equation}

The weights were set based on the relative scale of each loss term and held fixed across all experiments.

\section{Experiments}

\subsection{Experiment 1: Mechanism Comparison at 1.1M Parameters}

\paragraph{Setup.} We train TensionLM and an identical-architecture transformer on WikiText-2 for 10 epochs on CPU. The configuration is: $d_{\text{model}} = 128$, 4 layers, 4 heads, vocabulary size 2,048. The only difference between the two models is the attention mechanism: TensionLM uses sigmoid tension (Equations~\ref{eq:tau}--\ref{eq:output}) and the transformer uses standard softmax attention. All other hyperparameters, including learning rate, batch size, and sequence length, are identical.

\paragraph{Results.} Table~\ref{tab:exp1} shows the validation perplexity.

\begin{table}[h]
\centering
\caption{Mechanism comparison at 1.1M parameters on WikiText-2. Identical architecture and training; only the attention mechanism differs.}
\label{tab:exp1}
\vspace{4pt}
\begin{tabular}{llc}
\toprule
Model & Mechanism & Val PPL (WikiText-2) \\
\midrule
Transformer & Softmax attention & 57.8 \\
TensionLM & Sigmoid tension & \textbf{57.7} \\
\bottomrule
\end{tabular}
\end{table}

Sigmoid tension matches softmax attention exactly at this scale. The mechanism is viable --- the architectural change does not destroy training dynamics, degrade language modelling performance, or require additional tuning. This establishes the foundation: the theoretical motivation for switching from softmax to sigmoid does not come at a perplexity cost.

\subsection{Experiment 2: TS-Native Objectives at 13.5M Parameters}

\paragraph{Setup.} We train TensionLM (Small, 13.5M) on open-web-math \cite{paster2023openwebmath} for 1B tokens with and without the four TS-native auxiliary losses. Both models use the same architecture and cross-entropy objective; the TS-native variant adds ManifoldClosure (weight 0.01), TensionDiversity (0.02), ConstraintConsistency (0.1), and TensionEntropy (0.05).

\paragraph{Perplexity results.}

\begin{table}[h]
\centering
\caption{TS-native objectives at 13.5M parameters on open-web-math (1B tokens).}
\label{tab:exp2}
\vspace{4pt}
\begin{tabular}{lcc}
\toprule
Model & Objective & Val PPL \\
\midrule
Baseline & Cross-entropy only & 85.19 \\
TensionLM (TS-native) & CE + 4 aux losses & 86.50 \\
\bottomrule
\end{tabular}
\end{table}

The TS-native objectives cost 1.31 PPL relative to the baseline. This is the price of structural coherence in the constraint graph. We argue that this cost is acceptable and expected: the auxiliary losses constrain the model to produce representations that satisfy TS-predicted structural properties, which is a harder task than unconstrained next-token prediction. The question is not whether the cost exists, but what it buys.

\paragraph{Tension field analysis.} We inspect the $\tau$ field of both models on the prompt: \textit{``If A then B. If B then C. Therefore\ldots''} The TS-native model, Layer 5, Head 2, produces:

\begin{center}
\texttt{A[1]:0.32\quad B[3]:0.43\quad B[6]:0.64\quad C[8]:0.59\quad then[7]:0.50}
\end{center}

The full A$\rightarrow$B$\rightarrow$C transitivity chain is encoded as simultaneous active constraints. The tension values form a coherent chain: A influences B (0.32$\rightarrow$0.43), B influences C (0.64$\rightarrow$0.59), and the logical connective \textit{then} is held at moderate strength (0.50) as the structural bridge. The baseline model on the same prompt produces uniform diffuse activation with no identifiable chain structure --- the $\tau$ values are nearly equal across all positions, encoding no logical structure whatsoever.

\paragraph{Coherent text versus word salad.} We measure constraint graph statistics on coherent natural language text versus sequences of randomly ordered words from the same vocabulary. TS predicts that coherent text --- where constraints propagate and resolve --- should produce higher mean $\tau$ and denser active edge structure. The results confirm this: coherent text produces $+25\%$ mean $\tau$ and $+60\%$ active edge density compared to random word salad. The tension field measurably responds to coherence, as the theory predicts.

\subsection{Experiment 3: Curriculum Training at 13.5M Parameters}

\paragraph{Setup.} We train TensionLM (Small, 13.5M) sequentially on three domains:
\begin{enumerate}[label=(\arabic*)]
    \item \textbf{Logic}: Synthetic inference dataset comprising if/then chains, modus ponens, transitivity statements, and syllogisms formatted as natural language. 200M tokens. Final val PPL: 1.51.
    \item \textbf{Language}: WikiText-103. 100M tokens. Final val PPL: 196.
    \item \textbf{Mathematics}: open-web-math. 200M tokens. Final val PPL: 573.
\end{enumerate}
Each stage transfers the previous checkpoint with optimizer reset and learning rate divided by 3.

\paragraph{Motivation.} The curriculum design is driven by TS theory's prediction about formal data. General web text contains contradictory constraints: two documents asserting opposite facts about the same entity create constraint loops that never resolve. The model that trains on such data learns a fuzzy average --- a constraint graph in which most edges are weak and inconsistent because the training signal pushes them in conflicting directions simultaneously. Formal domains are different. A valid mathematical proof cannot contain a contradiction. A correct logical inference cannot have two incompatible conclusions. Training on formal data therefore forces the constraint graph to be coherent by construction.

We pre-load the constraint graph with pure logical structure in Stage 1, then expose it to natural language in Stage 2 (so the logical scaffold acquires vocabulary and syntax), and finally to mathematics in Stage 3.

\paragraph{First-contact perplexity ablation.}

\begin{table}[h]
\centering
\caption{First-contact mathematics perplexity (first validation checkpoint after switching to open-web-math). 13.5M scale ablation showing cumulative effect of curriculum stages.}
\label{tab:curriculum13m}
\vspace{4pt}
\begin{tabular}{lcc}
\toprule
Training History & Scale & First-Contact Maths PPL \\
\midrule
Cold start & 13.5M & $\approx$2293 \\
Logic only & 13.5M & $\approx$1076 \\
Logic + Language & 13.5M & $\approx$582 \\
\bottomrule
\end{tabular}
\end{table}

The logic-only curriculum halves the cold-start perplexity. Adding the language stage reduces it further to 582 --- a 4$\times$ improvement over cold start. Each stage provides cumulative benefit, consistent with the TS prediction that the constraint graph is progressively scaffolded.

\subsection{Experiment 4: 117M Curriculum Run}

\paragraph{Setup.} We scale the curriculum approach to 117M parameters using the Large configuration: $d_{\text{model}} = 768$, 12 layers, 12 heads, window $W = 64$ (local) with $W = 1024$ at every fourth global layer. Training uses 2$\times$ RTX 4090 GPUs with DDP, bf16 mixed precision, and \texttt{torch.compile}.

\begin{table}[h]
\centering
\caption{117M curriculum run: training stages, data, token budgets, and validation perplexity.}
\label{tab:exp4stages}
\vspace{4pt}
\begin{tabular}{llcc}
\toprule
Stage & Data & Tokens & Best Val PPL \\
\midrule
1 --- Logic & Synthetic inference & 200M & 5.10 \\
2 --- Language & FineWeb-Edu \cite{penedo2024fineweb} & 500M & 339 \\
3 --- Mathematics & open-web-math & 2B & 359.99 \\
\bottomrule
\end{tabular}
\end{table}

Stage 1 achieves val PPL 5.10 with train PPL $\approx$1.3, indicating the model has learned pure inference structure almost completely. Stage 2 uses FineWeb-Edu rather than WikiText-103: it is an educationally filtered subset of Common Crawl with substantially higher quality than general web text, more TS-aligned (fewer contradictory constraints), and provides 5$\times$ more tokens. Stage 3 trains for 2B tokens on open-web-math with minimum observed train PPL of 6.8.

For comparison, GPT-2 117M \cite{radford2019language} trained on general text achieves approximately 20 train PPL on comparable datasets. TensionLM achieves 6.8 on mathematics, a formally structured domain, having been pre-loaded with logical structure.

\paragraph{First-contact perplexity at 117M scale.}

\begin{table}[h]
\centering
\caption{First-contact mathematics perplexity at 117M scale. The curriculum effect amplifies strongly with model capacity.}
\label{tab:exp4firstcontact}
\vspace{4pt}
\begin{tabular}{lcc}
\toprule
Training History & Scale & First-Contact Maths PPL \\
\midrule
Cold start & 13.5M & $\approx$2293 \\
Logic + Language curriculum & 13.5M & $\approx$582 \\
\midrule
Logic + Language + Math curriculum & 117M & \textbf{24} \\
\bottomrule
\end{tabular}
\end{table}

The 117M curriculum model achieves first-contact train PPL of 24 on open-web-math, compared to approximately 2293 for an equivalent cold-start model --- a \textbf{96$\times$ improvement}. The curriculum effect does not simply scale linearly with parameter count: at 13.5M it produces 4$\times$ improvement; at 117M it produces 96$\times$. This superlinear amplification is consistent with the TS prediction that larger constraint graphs benefit disproportionately from pre-loading, because a larger graph can encode more specific logical scaffolding and therefore acquire new domains with less effort.

\paragraph{Stage 3 milestones.}

\begin{table}[h]
\centering
\caption{Key milestones during stage 3 (mathematics) training of the 117M model.}
\label{tab:milestones}
\vspace{4pt}
\begin{tabular}{cp{10cm}}
\toprule
Step & Event \\
\midrule
2,550 & First-contact train PPL: 24. No domain shock observed at curriculum transition. \\
2,550 & LaTeX notation appearing in generations. \\
2,550 & \texttt{To prove this by contradiction assume that $\theta_n \in [0,1]$} \\
10,750 & \texttt{The sum of angles in a triangle is 180\textdegree} --- correct. \\
10,750 & \texttt{lim(1/x) = 0} --- correct. \\
14,000 & Best validation checkpoint saved (val PPL 359.99). \\
29,900 & Val PPL 394 --- slight overfit in epoch 2, best checkpoint already saved. \\
\bottomrule
\end{tabular}
\end{table}

The absence of domain shock at curriculum transition is notable. First contact with mathematics at step 2,550 produces train PPL 24 immediately, without the large spike that would be expected if the model were encountering an entirely foreign distribution. This indicates successful transfer: the logical scaffold from stage 1, combined with the language competency from stage 2, provides a sufficiently coherent constraint graph to absorb mathematical notation with minimal disruption.

\subsection{Experiment 5: Formal Reasoning Evaluation}

\paragraph{Setup.} We evaluate the 117M curriculum model at the best checkpoint (step 14,000) on a hand-crafted 23-question formal reasoning benchmark spanning six categories: syllogisms, transitivity, arithmetic, calculus, algebra, and definitions. No instruction tuning, fine-tuning, or RLHF is applied. The model is evaluated as a raw pretrained language model.

\paragraph{Results.}

\begin{table}[h]
\centering
\caption{Formal reasoning evaluation on 23-question benchmark. 117M curriculum model, step 14,000 checkpoint, no fine-tuning.}
\label{tab:reasoning}
\vspace{4pt}
\begin{tabular}{lcc}
\toprule
Category & Questions & Accuracy \\
\midrule
Algebra & 3 & 67\% \\
Definitions & 3 & 67\% \\
Arithmetic & 4 & 50\% \\
Calculus & 4 & 50\% \\
Transitivity & 3 & 33\% \\
Syllogisms & 6 & 17\% \\
\midrule
\textbf{Overall} & \textbf{23} & \textbf{43.5\%} \\
\bottomrule
\end{tabular}
\end{table}

The model achieves 43.5\% overall on formal reasoning without any instruction tuning. Algebra and definitions score highest at 67\%, followed by arithmetic and calculus at 50\%. Transitivity (33\%) and syllogisms (17\%) score lowest.

\paragraph{Key finding: checkpoint selection.} The step 14,000 checkpoint outperforms the final checkpoint (step 30,518) on formal reasoning. The second epoch of mathematics training partially overwrites the logical inference structure accumulated during stage 1 of the curriculum. The model over-specialises on mathematical text at the expense of logical reasoning capacity. This is a direct TS prediction: dense new constraints of a different type can displace existing coherent constraint structure if training continues past an optimal point. The optimal reasoning checkpoint exists before full stage 3 saturation.

This finding has practical implications: for reasoning-critical applications, early stopping on formal reasoning benchmarks during domain-specific fine-tuning may preserve general logical capacity.

\section{Tension Field Analysis}

The tension field is the primary interpretability contribution of TensionLM. In this section we present qualitative and quantitative analysis of $\tau$ values from the 117M curriculum model.

\subsection{Simultaneous Non-Competitive Constraints}

The defining property of sigmoid tension over softmax is the ability to represent simultaneous full-strength constraint activation. We verify this on two prompts.

\paragraph{Named entity constraints.} For the prompt \textit{``Manchester United won the Premier League title''}, we inspect Layer 12, Head 0:

\begin{center}
\texttt{title $\leftarrow$ Manchester:0.95\quad United:0.95\quad League:0.86\quad Premier:0.71}
\end{center}

The token \textit{title} is simultaneously constrained by \textit{Manchester} and \textit{United} at $\tau = 0.95$ each. Under softmax, this is impossible: if two tokens each received 0.95 of the attention budget, the total would exceed 1.0. Under sigmoid tension, this is the natural and expected result: both tokens are genuine, independent constraints on the output, and both are active at full strength.

\paragraph{Logical syllogism constraints.} For the prompt \textit{``If all mammals are warm-blooded and all whales are mammals then''}, Layer 5, Head 2 produces:

\begin{center}
\texttt{then $\leftarrow$ mammals:0.755\quad whales:0.755\quad are:0.414}
\end{center}

Both logical subjects are held simultaneously at equal full strength (0.755 each). The word \textit{then} depends on both premises equally, and the constraint graph correctly encodes this dual dependency. Layer 5, Head 11 on the same prompt produces:

\begin{center}
\texttt{then $\leftarrow$ mammals:0.991}
\end{center}

Near-total constraint. Head 11 has specialised to identify the key inference term --- the concept that is most critical for determining what follows \textit{then} in a syllogistic context. The constraint graph knows what the output depends on before the output layer fires.

\subsection{Spontaneous Head Specialisation}

Without any explicit supervision over head roles, the 12 heads at 117M scale spontaneously differentiate into distinct functional types across training:

\begin{itemize}
    \item \textbf{Syntactic tracking heads}: high $\tau$ on grammatically related token pairs (subject-verb agreement, noun-adjective, determiner-noun). These heads respond strongly to local syntactic structure within phrases.
    \item \textbf{Long-range semantic heads}: high $\tau$ across sentence boundaries on topically related tokens. These heads track coreference, topic continuity, and thematic coherence across longer spans.
    \item \textbf{Diffuse background heads}: low, approximately uniform $\tau$ across the full window. These heads appear to provide a general context signal rather than specialised constraint tracking.
\end{itemize}

This specialisation emerges from the interaction of the tension diversity loss (which prevents heads collapsing onto a single position) and the constraint consistency loss (which enforces transitivity, implicitly rewarding heads that track consistent structural patterns). It is not supervised: no labels or auxiliary tasks specify which heads should perform which function.

\subsection{Summary of Tension Field Properties}

The tension field exhibits all properties predicted by TS theory: independent simultaneous constraint activation (impossible under softmax), sensitivity to logical structure (transitivity chains visible in $\tau$ values), sensitivity to coherence (+25\% mean $\tau$ and +60\% active edge density on coherent text versus word salad), and spontaneous specialisation into functionally distinct constraint tracking roles. The constraint graph is not a metaphor --- it is a directly readable structure produced by every forward pass.

\section{Discussion}

\subsection{Why the Curriculum Effect Amplifies with Scale}

The 96$\times$ improvement in first-contact mathematics perplexity at 117M, compared to 4$\times$ at 13.5M, is a superlinear scaling of the curriculum effect. This amplification is consistent with the TS interpretation. A larger model has a higher-capacity constraint graph. When this graph is pre-loaded with logical structure during stage 1, the scaffolding is more specific, more extensive, and more finely structured than the equivalent graph in a smaller model. When the model then encounters mathematics, the constraint graph has a richer set of pre-existing structures to connect new mathematical concepts to. The larger the graph, the more connections are available, and the less new structure needs to be built from scratch.

This predicts that the curriculum effect will continue to amplify at 1B and 7B parameter scale. We view this as a testable prediction of TS theory, to be verified in future work.

\subsection{The Optimal Reasoning Checkpoint and Constraint Displacement}

The finding that step 14,000 outperforms step 30,518 on formal reasoning is not a training instability --- it is a predicted consequence of TS dynamics. During stage 3, the model is exposed to 2B tokens of mathematical text. Mathematical constraints are dense and domain-specific. Over the second epoch, these constraints increasingly dominate the constraint graph, displacing the more general logical structure that was established in stage 1. The model's constraint graph becomes highly specialised for mathematical text at the expense of general logical inference capacity.

This suggests a practical training strategy: for models intended for general formal reasoning, stage 3 training should be stopped before epoch 2 completion on domain-specific data. The optimal point is determined by evaluating on a held-out formal reasoning benchmark during training and selecting the checkpoint that maximises reasoning accuracy rather than domain-specific perplexity.

\subsection{The Interpretability Case for TensionLM}

The mechanistic interpretability research programme has spent years developing techniques to reverse-engineer what transformer circuits compute: probing classifiers, activation patching, causal tracing, circuit analysis. These techniques exist because softmax transformers were not designed to be interpretable. Their attention patterns are normalised probability distributions, not constraint activations; their internal representations are not directly meaningful; understanding what they compute requires substantial post-hoc analysis.

TensionLM's $\tau$ tensor requires none of this. It is a directly readable constraint graph produced by every forward pass. A researcher who wants to know what the model thinks a particular output depends on can inspect the $\tau$ values for that output position and read them directly: token $A$ is constraining this output at strength 0.95, token $B$ at 0.755, and so on. The graph is the reasoning, not a shadow of the reasoning.

We view this as the primary contribution of TensionLM as an architecture, independent of performance metrics. The field has spent enormous effort trying to understand black-box models. TensionLM offers a glass box.

\subsection{Limitations and Future Work}

The present experiments are conducted at 117M parameters on two consumer GPUs. Competitive performance relative to production-scale language models requires substantially more compute: the standard scaling law literature \cite{hoffmann2022training} suggests that a 117M model trained on $\sim$2.7B tokens is approaching compute-optimal, but that meaningful capability improvements require scaling to 1B and beyond. We view 117M as a proof of concept.

The formal reasoning evaluation uses 23 hand-crafted questions. This is sufficient to identify qualitative patterns --- the advantage of algebra and calculus over syllogisms, the checkpoint selection effect --- but is too small to draw statistically robust conclusions. Evaluation on standardised benchmarks (GSM8K, MATH, LogiQA, ARC) is necessary for rigorous comparison with existing models.

The model presented here is a raw pretrained language model with no instruction tuning, no RLHF, and no chain-of-thought supervision. Applying SFT and RLHF to the best curriculum checkpoint is expected to substantially improve formal reasoning performance on standardised benchmarks.

Future work planned includes: Phase 4 ablations over window size, global layer frequency, tau-mass versus position-count normalisation, and RoPE versus learned embeddings; Phase 5 extending to longer context ($W = 512$ local, $W = 4096$ global); Phase 6 scaling to 1B and 7B parameters requiring compute grants; standardised benchmarks; and instruction fine-tuning.

\section{Conclusion}

We have presented TensionLM, a language model architecture grounded in Thinking System theory, whose central contribution is the replacement of softmax attention with sigmoid tension --- a change motivated not by efficiency but by theoretical correctness. Softmax encodes competition over a fixed budget. Real constraints are independent. Sigmoid tension implements independent constraint evaluation, producing a directly readable constraint graph as a byproduct of every forward pass.

At 1.1M parameters, the mechanism matches transformers exactly. At 13.5M with TS-native objectives, the constraint graph measurably encodes structural properties that cross-entropy alone does not produce: transitivity chains, sensitivity to coherence, and a 1.31 PPL cost that buys structurally coherent internal representations. At 117M with curriculum training, the model achieves first-contact mathematics perplexity 96$\times$ better than cold start --- a result that amplifies superlinearly with scale and predicts continued improvement at 1B and beyond.

The tension field is readable. \textit{Manchester} and \textit{United} both score 0.95 simultaneously --- impossible under softmax, natural under sigmoid tension. Both logical subjects in a syllogism score 0.755 simultaneously. Heads specialise without supervision. The constraint graph knows what the output depends on before the output layer fires. A model trained for 43.5\% formal reasoning accuracy with no instruction tuning demonstrates that curriculum pre-loading of logical structure provides substantial downstream reasoning capability.

Thinking System theory predicts that computation is constraint relaxation, that softmax encodes the wrong prior, and that formal training produces coherent constraint graphs. The experiments in this paper confirm all three predictions empirically. TensionLM is the glass box that mechanistic interpretability has been trying to reverse-engineer from black boxes. It is built in by design.

\bibliographystyle{plain}
\begin{thebibliography}{99}

\bibitem{vaswani2017attention}
A.~Vaswani, N.~Shazeer, N.~Parmar, J.~Uszkoreit, L.~Jones, A.~N.~Gomez, {\L}.~Kaiser, and I.~Polosukhin.
\newblock Attention is all you need.
\newblock In \textit{Advances in Neural Information Processing Systems}, volume 30, 2017.

\bibitem{radford2019language}
A.~Radford, J.~Wu, R.~Child, D.~Luan, D.~Amodei, and I.~Sutskever.
\newblock Language models are unsupervised multitask learners.
\newblock \textit{OpenAI Blog}, 1(8):9, 2019.

\bibitem{su2021roformer}
J.~Su, Y.~Lu, S.~Pan, A.~Murtadha, B.~Wen, and Y.~Liu.
\newblock {RoFormer}: Enhanced transformer with rotary position embedding.
\newblock \textit{arXiv preprint arXiv:2104.09864}, 2021.

\bibitem{shazeer2020glu}
N.~Shazeer.
\newblock {GLU} variants improve transformer.
\newblock \textit{arXiv preprint arXiv:2002.05202}, 2020.

\bibitem{bengio2009curriculum}
Y.~Bengio, J.~Louradour, R.~Collobert, and J.~Weston.
\newblock Curriculum learning.
\newblock In \textit{Proceedings of the 26th Annual International Conference on Machine Learning}, pages 41--48, 2009.

\bibitem{hoffmann2022training}
J.~Hoffmann, S.~Borgeaud, A.~Mensch, E.~Buchatskaya, T.~Cai, E.~Rutherford, D.~de~Las~Casas, L.~A.~Hendrycks, J.~Welbl, A.~Clark, et al.
\newblock Training compute-optimal large language models.
\newblock In \textit{Advances in Neural Information Processing Systems}, volume 35, 2022.

\bibitem{elhage2021mathematical}
N.~Elhage, N.~Nanda, C.~Olsson, T.~Henighan, N.~Joseph, B.~Mann, A.~Askell, Y.~Bai, A.~Chen, T.~Conerly, et al.
\newblock A mathematical framework for transformer circuits.
\newblock \textit{Transformer Circuits Thread}, 2021.

\bibitem{olsson2022context}
C.~Olsson, N.~Elhage, N.~Nanda, N.~Joseph, N.~DasSarma, T.~Henighan, B.~Mann, A.~Askell, Y.~Bai, A.~Chen, et al.
\newblock In-context learning and induction heads.
\newblock \textit{Transformer Circuits Thread}, 2022.

\bibitem{ramapuram2024sigmoid}
J.~Ramapuram, F.~Danieli, E.~Dhekane, S.~Webb, D.~Busbridge, P.~Ablin, A.~Korhonen, Y.~Zemlyanskiy, F.~Fleuret, and S.~Marks.
\newblock Theory, analysis, and best practices for sigmoid self-attention.
\newblock \textit{arXiv preprint arXiv:2409.04431}, 2024.

\bibitem{azerbayev2023llemma}
Z.~Azerbayev, H.~Schoelkopf, K.~Paster, M.~D.~Santos, S.~McAleer, A.~Q.~Jiang, J.~Deng, S.~Biderman, and S.~Welleck.
\newblock Llemma: An open language model for mathematics.
\newblock \textit{arXiv preprint arXiv:2310.01848}, 2023.

\bibitem{lewkowycz2022solving}
A.~Lewkowycz, A.~Andreassen, D.~Dohan, E.~Dyer, H.~Michalewski, V.~Ramasesh, A.~Slone, C.~Anil, I.~Schlag, T.~Gutman-Solo, et al.
\newblock Solving quantitative reasoning problems with language models.
\newblock In \textit{Advances in Neural Information Processing Systems}, volume 35, 2022.

\bibitem{hendrycks2021measuring}
D.~Hendrycks, C.~Burns, S.~Basart, A.~Zou, M.~Mazeika, D.~Song, and J.~Steinhardt.
\newblock Measuring massive multitask language understanding.
\newblock In \textit{International Conference on Learning Representations}, 2021.

\bibitem{cobbe2021training}
K.~Cobbe, V.~Kosaraju, M.~Bavarian, M.~Chen, H.~Jun, L.~Kaiser, M.~Plappert, J.~Tworek, J.~Hilton, R.~Nakano, et al.
\newblock Training verifiers to solve math word problems.
\newblock \textit{arXiv preprint arXiv:2110.14168}, 2021.

\bibitem{wei2022chain}
J.~Wei, X.~Wang, D.~Schuurmans, M.~Bosma, B.~Ichter, F.~Xia, E.~Chi, Q.~Le, and D.~Zhou.
\newblock Chain-of-thought prompting elicits reasoning in large language models.
\newblock In \textit{Advances in Neural Information Processing Systems}, volume 35, 2022.

\bibitem{hopfield1982neural}
J.~J.~Hopfield.
\newblock Neural networks and physical systems with emergent collective computational abilities.
\newblock \textit{Proceedings of the National Academy of Sciences}, 79(8):2554--2558, 1982.

\bibitem{lecun2006tutorial}
Y.~LeCun, S.~Chopra, R.~Hadsell, M.~Ranzato, and F.~Huang.
\newblock A tutorial on energy-based learning.
\newblock In G.~Bakir, T.~Hofman, B.~Scholkopf, A.~Smola, and B.~Taskar, editors, \textit{Predicting Structured Data}. MIT Press, 2006.

\bibitem{friston2010free}
K.~Friston.
\newblock The free-energy principle: A unified brain theory?
\newblock \textit{Nature Reviews Neuroscience}, 11(2):127--138, 2010.

\bibitem{penedo2024fineweb}
G.~Penedo, H.~Kydl{\'\i}{\v{c}}ek, L.~allal, A.~Lozhkov, M.~Mitchell, C.~Raffel, L.~Von~Werra, and T.~Wolf.
\newblock The {FineWeb} datasets: Decanting the web for the finest text data at scale.
\newblock \textit{arXiv preprint arXiv:2406.17557}, 2024.

\bibitem{paster2023openwebmath}
K.~Paster, M.~D.~Santos, Z.~Azerbayev, and J.~Ba.
\newblock {OpenWebMath}: An open dataset of high-quality mathematical web text.
\newblock \textit{arXiv preprint arXiv:2310.06786}, 2023.

\bibitem{sennrich2016neural}
R.~Sennrich, B.~Haddow, and A.~Birch.
\newblock Neural machine translation of rare words with subword units.
\newblock In \textit{Proceedings of the 54th Annual Meeting of the Association for Computational Linguistics}, pages 1715--1725, 2016.

\bibitem{xiao2023efficient}
G.~Xiao, Y.~Tian, B.~Chen, S.~Han, and M.~Lewis.
\newblock Efficient streaming language models with attention sinks.
\newblock \textit{arXiv preprint arXiv:2309.17453}, 2023.

\end{thebibliography}

\end{document}
